%%=============================================================================
%% Voorwoord
%%=============================================================================

\chapter*{\IfLanguageName{dutch}{Woord vooraf}{Preface}}%
\label{ch:voorwoord}

%% TODO:
%% Het voorwoord is het enige deel van de bachelorproef waar je vanuit je
%% eigen standpunt (``ik-vorm'') mag schrijven. Je kan hier bv. motiveren
%% waarom jij het onderwerp wil bespreken.
%% Vergeet ook niet te bedanken wie je geholpen/gesteund/... heeft

% \lipsum[1-2]

De kern van deze bachelorproef is ontstaan uit een idee dat eerst werd voorgesteld door Thomas Clauwaert -- docent voor (onder andere) het \emph{OLOD} (opleidingsonderdeel) \emph{Cybersecurity Advanced} binnen de opleiding \emph{Toegepaste Informatica} op \emph{HOGENT}. Zijn initiële vraag -- hoe een infrastructuur, gevisualiseerd in een graaf, geautomatiseerd kan worden opgezet door \emph{configuration management} tools -- heb ik vanuit een nieuwe lens bekeken en uitgewerkt tot de uiteindelijke tool die in dit werk wordt beschreven.

Dit proces was bijzonder uitdagend, en het pad naar het einde was allesbehalve rechtlijnig. Het bouwen van de brug tussen punt A (de diagrammen) en punt B (de effectieve omgeving) heeft veel tijd, energie en doorzettingsvermogen van mij gevraagd. Wat me heeft doen doorzetten, is het keer op keer oplossen van het ene probleem na het andere -- tot ik uiteindelijk een resultaat bereikt heb waar ik oprecht trots op ben.
Ik ben er sterk van overtuigd dat ik mij nog zal blijven bezighouden met dit project (ook na mijn studies), in de hoop dat dit de drempel zal verlagen voor docenten en studenten om kwalitatieve IT-omgevingen sneller en efficiënter te kunnen opzetten.

Graag wil ik hier de kans nemen om een aantal mensen bedanken.

Allereerst mijn promotor, \textbf{Alexander Veldeman}, voor zijn begeleiding, gerichte feedback en steun doorheen beide semesters, alsook voor de oprechte interesse die hij heeft getoond in mijn idee en de tool die eruit is ontstaan.

Vervolgens mijn co-promotor, \textbf{Yoeri Rousseaux} van District09 (mijn stagebedrijf), die mij op technisch vlak meermaals heeft geholpen om het noorden terug te vinden dankzij zijn lange werkervaring in de IT-sector.

Uiteraard gaat een bijzonder dankwoord uit naar \textbf{Thomas Clauwaert}, voor het initiële idee dat aan de basis ligt van deze bachelorproef, en voor zijn boeiende lessen die ik heb kunnen volgen tijdens \emph{Cybersecurity \& Virtualisation} en \emph{Cybersecurity Advanced}.

Daarnaast wil ik \textbf{Bert Van Vreckem} bedanken voor zijn bijdrage aan de opleiding en, bij uitbreiding, aan deze bachelorproef. Zonder zijn \emph{Ansible}-rollen en \emph{Ansible Skeleton} zou deze tool niet de vorm hebben kunnen aannemen die deze vandaag heeft.

Ten slotte wil ik ook de andere docenten van de opleiding \emph{Toegepaste Informatica} bedanken voor de kennis, ervaring en passie die zij me hebben meegegeven tijdens mijn jaren op \emph{HOGENT}.

\begin{flushright}
    \emph{Lucas Ludueña-Segre, Gent, mei 2026}
\end{flushright}